\documentclass[a4paper,12pt]{article}
\usepackage[utf8]{inputenc}
\usepackage[T1]{fontenc}
\usepackage[german]{babel}
\usepackage{amsmath}
\usepackage{amssymb}
\usepackage{amsthm}
\usepackage{geometry}
\usepackage{graphicx}
\usepackage[hidelinks]{hyperref}
\usepackage{listings}
\usepackage{xcolor}
\usepackage{enumitem}
\usepackage{booktabs}
\usepackage{tikz}

\geometry{margin=2.5cm}

\theoremstyle{definition}
\newtheorem{definition}{Definition}
\newtheorem{beispiel}{Beispiel}
\newtheorem{satz}{Satz}
\newtheorem{lemma}{Lemma}
\newtheorem{korollar}{Korollar}
\newtheorem{bemerkung}{Bemerkung}

\setlist[itemize,1]{label=--}

\DeclareMathOperator{\popcnt}{popcnt}
\newcommand{\code}[1]{\texttt{\detokenize{#1}}}

% ---------- Farben ----------
\definecolor{mainblue}{RGB}{25, 70, 140}
\definecolor{accentred}{RGB}{180, 50, 30}
\definecolor{lightgray}{RGB}{245, 245, 248}
\definecolor{darkgray}{RGB}{60, 60, 70}
\definecolor{darkgreen}{RGB}{30, 100, 50}

\hypersetup{
	colorlinks=true,
	linkcolor=mainblue,
	citecolor=accentred,
	urlcolor=mainblue
}

%  Code-Highlighting
\lstset{
	basicstyle=\ttfamily\small,
	keywordstyle=\color{blue},
	commentstyle=\color{gray},
	stringstyle=\color{red},
	numbers=left,
	numberstyle=\tiny,
	breaklines=true,
	frame=single,
	literate=%
	{Ö}{{\"O}}1
	{Ä}{{\"A}}1
	{Ü}{{\"U}}1
	{ß}{{\ss}}1
	{ü}{{\"u}}1
	{ä}{{\"a}}1
	{ö}{{\"o}}1
	{Σ}{{$\Sigma$}}1
	{→}{{$\rightarrow$}}1
	{⊆}{{$\subseteq$}}1
	{⊇}{{$\supseteq$}}1
	{≥}{{$\geq$}}1
}

\title{Der Subnetting Algorithmus sbntgg}
\author{Stephan Epp}
\date{18. Oktober 2026}

\begin{document}

\maketitle

\begin{abstract}
\noindent
Der Subnetting Algorithmus \texttt{sbntgg} (Kommandozeilenprogramm \texttt{subnetplan}) teilt einen zugewiesenen IPv4- oder IPv6-Adressbereich automatisch in überlappungsfreie Subnetze für Mitarbeiter- und Ressourcengruppen auf und erzeugt aus einer Zugriffsmatrix Firewall-Regeln für \texttt{iptables}, \texttt{nftables} und Cisco-ACLs.
Die Arbeit definiert das Problem formal, beweist die Korrektheit der Zuteilung (Ausrichtung, Überlappungsfreiheit, Kapazitätsgarantie), zeigt, dass die Strategie \glqq größte Blöcke zuerst\grqq{} lückenlos und unter allen Anordnungen optimal ist, charakterisiert die freien Restbereiche exakt und beweist, dass sich jede Wachstumsreserve ohne Umnummerierung nutzen lässt, sobald ihre Gesamtgröße in den freien Adressraum passt.
Für die Firewall-Regeln wird ein Default-Deny-Satz bewiesen.
Die Laufzeit des Kerns beträgt $\Theta(n\log n + a)$ im Vergleichsmodell und ist damit optimal.
Ein ausführlich erklärtes Konfigurationsbeispiel, eine Einführung in das Subnetting für Schülerinnen und Schüler sowie ein umfangreicher Ausblick ergänzen die Arbeit.
\end{abstract}

\tableofcontents
\newpage

\section{Einführung}

Der Subnetting Algorithmus \texttt{sbntgg} ist ein deterministischer Algorithmus zur automatischen Planung von IP-Netzen. Aus einer Konfigurationsdatei mit einem zugewiesenen Adressbereich (dem \emph{Basisnetz}) und einer Liste von Gruppen (zum Beispiel \glqq Entwicklung\grqq, \glqq HR\grqq, \glqq Drucker\grqq) berechnet er für jede Gruppe ein eigenes Subnetz mit VLAN-ID, Gateway und Wachstumsreserve. Optional erzeugt er aus der Zugriffsmatrix der Gruppen ein Regelwerk für Firewalls.

\subsection{Problemstellung}

Gegeben sind ein Basisnetz $B$ (IPv4 oder IPv6), $n$ Gruppen $g_1,\ldots,g_n$ mit jeweils $h_i$ Hosts, ein Wachstumsfaktor $\gamma \ge 1$, eine Strategie (\code{vlsm} oder \code{fixed}), eine Gateway-Position (\code{first}, \code{last} oder \code{none}) sowie pro Gruppe optional eine VLAN-ID, eine Priorität und eine Menge von Gruppen, auf die sie zugreifen darf.

\textbf{Ziel:} Bestimme für jede Gruppe ein Subnetz innerhalb von $B$, sodass alle Subnetze paarweise überlappungsfrei sind, jedes Subnetz mindestens $\lceil \gamma h_i \rceil$ nutzbare Host-Adressen besitzt und jede Gruppe eine eindeutige VLAN-ID erhält.

\textbf{Ergebnis:}
\begin{itemize}
	\item Wenn die Summe der benötigten Blockgrößen in $B$ passt: Gib den Plan aus (Subnetze, Gateways, Host-Bereiche, freie Restbereiche) und, falls gewünscht, Firewall-Regeln.
	\item Wenn die Summe nicht in $B$ passt: Brich mit einer Fehlermeldung zum Platzmangel ab und gib keinen Plan aus (Satz~\ref{satz:existenz} zeigt, dass dann auch jedes andere Verfahren scheitern muss).
	\item Wenn die Konfiguration ungültig ist (unbekannte Schlüssel, doppelte Namen oder VLAN-IDs, unbekannte Zugriffsziele): Brich ebenfalls mit einer Fehlermeldung ab.
	\item Wenn \code{access} leer ist: Die Gruppe erhält keinen Zugriff auf interne Netze.
\end{itemize}

\subsection{Motivation}

Die manuelle Planung von Subnetzen ist fehleranfällig: Netze überlappen sich, Blockgrenzen werden falsch berechnet, VLAN-IDs doppelt vergeben, und das Wachstum der Gruppen wird vergessen. Gleichzeitig wird die Zugriffsmatrix \glqq wer darf auf wen\grqq{} meist getrennt vom Adressplan in Firewall-Regeln übersetzt, wodurch Plan und Regeln auseinanderlaufen. \texttt{sbntgg} berechnet beides aus einer einzigen Quelle.

Eine besondere Verbindung besteht zum Subgraph Algorithmus: Die Zugriffsmatrix ist die Adjazenzmatrix eines gerichteten Graphen über den Gruppen (Abschnitt~\ref{sec:subgraph-verbindung}). Der Vergleich zweier Richtlinien ist damit ein Subgraph-Problem.

\subsection{Aufbau der Arbeit}

Abschnitt~2 erklärt das Subnetting zuerst anschaulich und dann formal. Abschnitt~3 beschreibt den Algorithmus samt Implementierung und beweist seine Korrektheit. Abschnitt~4 analysiert Laufzeit, Optimalität und Speicher. Abschnitt~5 erläutert das Konfigurationsbeispiel und weitere Code-Beispiele. Abschnitt~6 fasst zusammen und gibt einen ausführlichen Ausblick.

\section{Grundlagen}

\subsection{Subnetting leicht erklärt}
\label{sec:leicht}

Dieser Abschnitt ist so geschrieben, dass er auch in einer Realschulklasse (etwa im Informatikunterricht ab Klasse~8) verstanden werden kann. Es wird nur Rechnen mit Zweierpotenzen benötigt.

\subsubsection{Was ist eine IP-Adresse?}

Jedes Gerät in einem Netzwerk (Computer, Handy, Drucker) braucht eine Adresse, damit Nachrichten ankommen, so wie ein Haus eine Hausnummer braucht, damit der Briefträger es findet. Eine IPv4-Adresse besteht aus vier Zahlen zwischen 0 und 255, getrennt durch Punkte, zum Beispiel \texttt{10.20.1.130}.

Warum nur bis 255? Jede der vier Zahlen besteht aus 8 \emph{Bit}. Ein Bit ist eine Stelle, die 0 oder 1 sein kann. Mit 8 Bit gibt es $2^8 = 256$ Möglichkeiten, also die Zahlen 0 bis 255. Eine ganze Adresse hat $4\cdot 8 = 32$ Bit und damit $2^{32} = 4\,294\,967\,296$ mögliche Adressen.

\subsubsection{Das Dualsystem in einer Minute}

Computer rechnen mit Zweierpotenzen. Jede Stelle einer Dualzahl hat einen Wert, der doppelt so groß ist wie die Stelle rechts davon:
\begin{center}
\begin{tabular}{lcccccccc}
\toprule
Stellenwert & 128 & 64 & 32 & 16 & 8 & 4 & 2 & 1 \\
\midrule
$10$ & 0 & 0 & 0 & 0 & 1 & 0 & 1 & 0 \\
$130$ & 1 & 0 & 0 & 0 & 0 & 0 & 1 & 0 \\
$192$ & 1 & 1 & 0 & 0 & 0 & 0 & 0 & 0 \\
$224$ & 1 & 1 & 1 & 0 & 0 & 0 & 0 & 0 \\
\bottomrule
\end{tabular}
\end{center}
Zum Beispiel ist $130 = 128 + 2$, also steht eine 1 bei 128 und eine 1 bei 2.

\subsubsection{Netz und Host: Straße und Hausnummer}

Eine Adresse besteht aus zwei Teilen: einem \emph{Netzteil} (die Straße) und einem \emph{Hostteil} (die Hausnummer in dieser Straße). Der Zusatz hinter dem Schrägstrich, etwa \texttt{/26}, sagt, wie viele der 32 Bit zum Netzteil gehören. Das ist die \emph{Präfixlänge}. Die übrigen Bit sind der Hostteil.

Beim Netz \texttt{10.20.1.128/26} gehören 26 Bit zur Straße. Für die Hausnummern bleiben $32-26 = 6$ Bit, also $2^6 = 64$ Adressen. Wichtig ist: \textbf{Je größer die Zahl hinter dem Schrägstrich, desto kleiner das Netz.}

\begin{center}
\begin{tabular}{rrrr}
\toprule
Präfix & Hostbits & Adressen im Netz & Subnetzmaske \\
\midrule
\texttt{/24} & 8 & 256 & \texttt{255.255.255.0} \\
\texttt{/25} & 7 & 128 & \texttt{255.255.255.128} \\
\texttt{/26} & 6 & 64 & \texttt{255.255.255.192} \\
\texttt{/27} & 5 & 32 & \texttt{255.255.255.224} \\
\texttt{/28} & 4 & 16 & \texttt{255.255.255.240} \\
\texttt{/29} & 3 & 8 & \texttt{255.255.255.248} \\
\texttt{/30} & 2 & 4 & \texttt{255.255.255.252} \\
\bottomrule
\end{tabular}
\end{center}

Die \emph{Subnetzmaske} ist dieselbe Information anders geschrieben: so viele Einsen von links, wie das Präfix lang ist. Bei \texttt{/26} sind das 26 Einsen, im letzten Byte \texttt{11000000}, das ist $128 + 64 = 192$.

\subsubsection{Drei besondere Adressen}

In jedem Netz sind einige Adressen nicht für normale Geräte frei:
\begin{itemize}
	\item Die \textbf{Netzadresse} (die erste Adresse) ist das Straßenschild. Sie benennt das ganze Netz.
	\item Die \textbf{Broadcast-Adresse} (die letzte Adresse) ist die Lautsprecherdurchsage \glqq an alle in dieser Straße\grqq.
	\item Das \textbf{Gateway} ist das Tor aus der Straße hinaus in andere Straßen und ins Internet. Man legt es üblicherweise auf die erste oder letzte freie Adresse.
\end{itemize}
Von den Adressen eines Netzes gehen also bis zu drei für Geräte verloren. Ein \texttt{/28} hat 16 Adressen, davon sind $16 - 3 = 13$ für Geräte nutzbar.

\subsubsection{Das Rezept: Wie groß muss das Netz sein?}

Beispiel: In der Gruppe \glqq Drucker\grqq{} stehen 12 Geräte, und es soll 25\,\% Platz für Wachstum geben (Faktor $1{,}25$).
\begin{enumerate}
	\item Geräte: $12$.
	\item Mit Wachstum: $12 \cdot 1{,}25 = 15$ (bei Kommazahlen wird aufgerundet).
	\item Plus drei Sonderadressen: $15 + 3 = 18$.
	\item Die nächste Zweierpotenz, die mindestens 18 ist: $2^5 = 32$ (denn $2^4 = 16$ ist zu klein).
	\item Hostbits sind $5$, also ist das Präfix $32 - 5 = 27$. Das Netz ist ein \texttt{/27}.
\end{enumerate}
Probe: Ein \texttt{/27} hat 32 Adressen, davon sind $32-3 = 29$ nutzbar. Das ist mehr als 15. Ein \texttt{/28} hätte nur 13 nutzbare Adressen gehabt, zu wenig.

\subsubsection{Die Schokoladentafel: Warum Zweierpotenzen?}

Stell dir eine Schokoladentafel vor, die nur in der Mitte gebrochen werden kann. Jede Hälfte hat wieder eine Mitte, und so fort. So entstehen immer Stücke mit $1, 2, 4, 8, \dots$ Teilen, niemals Stücke mit 3 oder 5 Teilen. Genau so werden Adressbereiche geteilt (Abbildung~\ref{fig:halbierung}): Jede Halbierung erhöht das Präfix um 1.

\begin{figure}[ht]
\centering
\begin{tikzpicture}[x=1cm,y=1cm,
	blk/.style={draw,minimum height=0.7cm,inner sep=0pt,font=\footnotesize,fill=lightgray}]
	% Ebene /24
	\node[blk,minimum width=12.8cm,fill=mainblue!15] at (6.4,3.2) {\texttt{/24}: 256 Adressen};
	% Ebene /25
	\node[blk,minimum width=6.4cm,fill=mainblue!10] at (3.2,2.2) {\texttt{/25}: 128};
	\node[blk,minimum width=6.4cm,fill=mainblue!10] at (9.6,2.2) {\texttt{/25}: 128};
	% Ebene /26
	\foreach \i in {0,1,2,3}{
		\node[blk,minimum width=3.2cm,fill=mainblue!6] at (1.6+3.2*\i,1.2) {\texttt{/26}: 64};
	}
	% Ebene /27
	\foreach \i in {0,...,7}{
		\node[blk,minimum width=1.6cm,fill=white] at (0.8+1.6*\i,0.2) {\texttt{/27}: 32};
	}
\end{tikzpicture}
\caption{Jede Halbierung verdoppelt die Anzahl der Blöcke und halbiert ihre Größe. Ein Block beginnt immer an einer \glqq Rille\grqq.}
\label{fig:halbierung}
\end{figure}

Daraus folgen zwei Regeln, die der Algorithmus befolgt:
\begin{enumerate}
	\item Ein Netz darf nur dort beginnen, wo eine \glqq Rille\grqq{} ist (\emph{Ausrichtung}). Ein \texttt{/27} mit 32 Adressen kann bei den Adressen $0, 32, 64, \dots$ beginnen, aber nicht bei 40.
	\item Man legt \textbf{die großen Stücke zuerst} und die kleinen danach. Dann bleiben keine Lücken, weil jedes kleinere Stück genau dort beginnen kann, wo das vorige aufhört.
\end{enumerate}

\begin{beispiel}[Eine Schule mit \texttt{192.168.0.0/24}]
	Die Schule hat 256 Adressen. Es gibt 100 Schüler-PCs, 25 Lehrer-PCs und 10 Verwaltungs-PCs, ohne Wachstumsreserve und mit Gateway. Das Programm rechnet:
	\begin{center}
	\begin{tabular}{lrrrrl}
	\toprule
	Gruppe & Geräte & $+3$ & Zweierpotenz & Präfix & Netz \\
	\midrule
	Schüler & 100 & 103 & $128 = 2^7$ & \texttt{/25} & \texttt{192.168.0.0/25} \\
	Lehrer & 25 & 28 & $32 = 2^5$ & \texttt{/27} & \texttt{192.168.0.128/27} \\
	Verwaltung & 10 & 13 & $16 = 2^4$ & \texttt{/28} & \texttt{192.168.0.160/28} \\
	\bottomrule
	\end{tabular}
	\end{center}
	Die Schüler (das größte Netz) kommen zuerst: Adressen 0 bis 127. Danach folgen die Lehrer ab 128 bis 159 und die Verwaltung ab 160 bis 175. Belegt sind $128 + 32 + 16 = 176$ von 256 Adressen ($68{,}75\,\%$), frei sind die 80 Adressen \texttt{192.168.0.176/28} und \texttt{192.168.0.192/26}.
\end{beispiel}

\subsubsection{Wozu das Ganze?}

\begin{itemize}
	\item \textbf{Ordnung:} Jede Abteilung hat ihren eigenen Bereich, und man erkennt an der Adresse, wohin ein Gerät gehört.
	\item \textbf{Sicherheit:} Gäste-WLAN und Server liegen in verschiedenen Netzen. Eine Firewall entscheidet, wer von wo nach wo darf.
	\item \textbf{Geschwindigkeit:} Eine Durchsage (Broadcast) erreicht nur die Geräte im eigenen kleinen Netz und nicht den ganzen Betrieb.
	\item \textbf{Wachstum:} Freie Bereiche bleiben für neue Gruppen übrig.
\end{itemize}

\paragraph{Kontrollfragen.} (1) Wie viele Adressen hat ein \texttt{/29}? ($2^3 = 8$.) (2) Wie viele davon sind für Geräte nutzbar, wenn es Netzadresse, Broadcast und Gateway gibt? ($8-3=5$.) (3) Welches Präfix braucht man für 50 Geräte ohne Reserve? ($50+3 = 53 \le 64 = 2^6$, also \texttt{/26}.)

\subsection{Formale Definitionen}

Im Folgenden bezeichnet $w \in \{32, 128\}$ die Adressbreite (IPv4 beziehungsweise IPv6). Adressen werden als ganze Zahlen identifiziert.

\begin{definition}[Adressraum und Block]
	Der Adressraum ist $\mathcal{A}_w = \{0, 1, \ldots, 2^w - 1\}$. Für eine Präfixlänge $p \in \{0,\ldots,w\}$ und eine Adresse $a \in \mathcal{A}_w$ mit $2^{w-p} \mid a$ heißt
	\[
	N(a,p) = \{a, a+1, \ldots, a + 2^{w-p} - 1\}
	\]
	ein \emph{Block} (CIDR-Block) der Präfixlänge $p$. Er hat die Größe $|N(a,p)| = 2^{w-p}$, die \emph{Netzadresse} $a$ und die \emph{letzte Adresse} $a + 2^{w-p} - 1$. Die Bedingung $2^{w-p} \mid a$ heißt \emph{Ausrichtung}.
\end{definition}

\begin{lemma}[Laminarität]
\label{lem:laminar}
	Zwei Blöcke $N_1 = N(a_1,p_1)$ und $N_2 = N(a_2,p_2)$ mit $p_1 \ge p_2$ sind entweder disjunkt oder es gilt $N_1 \subseteq N_2$. Bei $p_1 = p_2$ sind sie gleich oder disjunkt.
\end{lemma}

\begin{proof}
	Sei $q_i = w - p_i$, also $q_1 \le q_2$. Wegen der Ausrichtung ist $N_i = \{x : \lfloor x / 2^{q_i} \rfloor = a_i / 2^{q_i}\}$. Gilt $N_1 \cap N_2 \neq \emptyset$, so gibt es ein $x$ in beiden Mengen. Für jedes $y \in N_1$ gilt $\lfloor y/2^{q_1}\rfloor = a_1/2^{q_1} = \lfloor x/2^{q_1}\rfloor$. Mit der Identität $\lfloor \lfloor z/m \rfloor / k \rfloor = \lfloor z/(mk) \rfloor$ für natürliche $m,k$ folgt
	\[
	\left\lfloor \frac{y}{2^{q_2}} \right\rfloor = \left\lfloor \frac{\lfloor y/2^{q_1}\rfloor}{2^{q_2-q_1}} \right\rfloor = \left\lfloor \frac{\lfloor x/2^{q_1}\rfloor}{2^{q_2-q_1}} \right\rfloor = \left\lfloor \frac{x}{2^{q_2}} \right\rfloor = \frac{a_2}{2^{q_2}},
	\]
	also $y \in N_2$. Somit $N_1 \subseteq N_2$. Bei $p_1 = p_2$ gilt zusätzlich $N_2 \subseteq N_1$ mit demselben Argument, also $N_1 = N_2$.
\end{proof}

\begin{definition}[Gruppe, Konfiguration]
	Eine \emph{Gruppe} ist ein Tupel $g = (\text{name}, h, v, \pi, \text{acc})$ mit Namen, Hostzahl $h \in \mathbb{N}_{\ge 1}$, optionaler VLAN-ID $v \in \{1,\ldots,4094\}$, Priorität $\pi \in \mathbb{N}_0$ (Standard $0$) und einer Liste $\text{acc}$ von Gruppennamen (Standard: leer). Eine \emph{Konfiguration} besteht aus einem Basisnetz $B = N(a_B, p_B)$, einer Liste $g_1,\ldots,g_n$ ($n \ge 1$) von Gruppen mit paarweise verschiedenen Namen, einem Wachstumsfaktor $\gamma \ge 1$, einer Strategie $\in \{\text{vlsm}, \text{fixed}\}$ und einer Gateway-Position $\in \{\text{first}, \text{last}, \text{none}\}$. Die Reihenfolge der Liste ist die \emph{Dateireihenfolge}.
\end{definition}

\begin{definition}[Bedarf, Blockexponent]
\label{def:bedarf}
	Der \emph{Overhead} ist $o = 2 + [\text{gw} \ne \text{none}]$ für IPv4 und $o = 1 + [\text{gw} \ne \text{none}]$ für IPv6, wobei $[\cdot]$ die Iverson-Klammer ist (IPv6 kennt keine Broadcast-Adresse). Der \emph{Bedarf} der Gruppe $g_i$ ist
	\[
	r_i = \lceil \gamma h_i \rceil + o.
	\]
	Der \emph{Blockexponent} ist $k_i = \lceil \log_2 r_i \rceil$, die \emph{Blockgröße} $s_i = 2^{k_i}$ und die \emph{Präfixlänge} $p_i = w - k_i$. Ferner sei $S = \sum_{i=1}^{n} s_i$.
\end{definition}

\begin{definition}[Subnetz]
\label{def:subnetz}
	Ein \emph{Subnetz} der Gruppe $g_i$ ist ein Block $N = N(a,p)$ mit $s = 2^{w-p}$ zusammen mit folgenden ausgezeichneten Adressen (Netzadresse $a$ und, nur IPv4, Broadcast $a+s-1$):
	\begin{center}
	\begin{tabular}{lcccc}
	\toprule
	Gateway-Position & Gateway & erster Host & letzter Host (IPv4) & letzter Host (IPv6) \\
	\midrule
	\code{none}  & -- & $a+1$ & $a+s-2$ & $a+s-1$ \\
	\code{first} & $a+1$ & $a+2$ & $a+s-2$ & $a+s-1$ \\
	\code{last}  & $a+s-2$ (IPv4), $a+s-1$ (IPv6) & $a+1$ & $a+s-3$ & $a+s-2$ \\
	\bottomrule
	\end{tabular}
	\end{center}
	Die Anzahl \emph{nutzbarer Hosts} ist $u = (\text{letzter Host}) - (\text{erster Host}) + 1$.
\end{definition}

\begin{lemma}[Nutzbare Hosts]
\label{lem:usable}
	Für jedes Subnetz der Größe $s$ gilt $u = s - o$. Ist $s \ge r_i$, so gilt $u \ge \lceil \gamma h_i \rceil \ge 1$.
\end{lemma}

\begin{proof}
	IPv4 ohne Gateway: $u = (a+s-2) - (a+1) + 1 = s - 2 = s - o$. IPv4 mit Gateway (\code{first}): $u = (a+s-2) - (a+2) + 1 = s-3 = s - o$; (\code{last}): $u = (a+s-3) - (a+1) + 1 = s - 3 = s-o$. IPv6 ohne Gateway: $u = (a+s-1) - (a+1) + 1 = s-1 = s-o$; mit Gateway: $u = s-2 = s-o$ in beiden Positionen. Ist $s \ge r_i = \lceil\gamma h_i\rceil + o$, so folgt $u = s - o \ge \lceil\gamma h_i\rceil$, und $\lceil \gamma h_i \rceil \ge 1$ wegen $\gamma h_i \ge h_i \ge 1$.
\end{proof}

\begin{definition}[Plan]
	Ein \emph{Plan} ist eine Liste von Subnetzen $(N_1,\ldots,N_n)$ (in Adressreihenfolge ausgegeben) mit zugehörigen VLAN-IDs, zusammen mit der Liste der \emph{freien Bereiche}, einer Zerlegung von $B \setminus \bigcup_i N_i$ in Blöcke. Die \emph{belegten Adressen} sind $\sum_i |N_i|$, die \emph{freien Adressen} sind $F = |B| - \sum_i |N_i|$ und die \emph{Auslastung} ist $100 \cdot \sum_i |N_i| / |B|$ Prozent.
\end{definition}

\begin{definition}[Zugriffsmatrix und Zugriffsgraph]
\label{def:zugriff}
	Die \emph{Zugriffsmatrix} ist die Matrix $Z \in \{0,1\}^{n\times n}$ mit $Z_{ij} = 1$ genau dann, wenn der Name von $g_j$ in $\text{acc}(g_i)$ vorkommt. Der \emph{Zugriffsgraph} ist der gerichtete Graph $G_Z = (V,E)$ mit $V = \{g_1,\ldots,g_n\}$ und $E = \{(g_i,g_j) : Z_{ij} = 1\}$ (vgl.\ die Definition der Adjazenzmatrix im Subgraph Algorithmus). Die Menge der \emph{Zugriffspaare} ist $P = \{(i,j) : Z_{ij} = 1,\ i \neq j\}$.
\end{definition}

\subsection{Grundidee}

Der Algorithmus basiert auf drei zentralen Konzepten.

\subsubsection{Blockgrößen als Zweierpotenzen}

Jede Gruppe erhält den kleinsten Block, der ihren Bedarf fasst. Blöcke sind nach Lemma~\ref{lem:laminar} entweder disjunkt oder ineinander enthalten. Das macht die Überlappungsfreiheit zu einer reinen Frage der Startadressen.

\subsubsection{Größte Blöcke zuerst (VLSM)}

\emph{Variable Length Subnet Masks} (VLSM) bedeutet, dass Subnetze unterschiedlich groß sein dürfen. Werden die Blöcke nach absteigender Größe hintereinander gelegt, ist jede Startadresse automatisch ausgerichtet (Satz~\ref{satz:luecke}). Es entstehen keine Lücken und der Platzbedarf ist minimal (Satz~\ref{satz:kompakt}).

\begin{beispiel}
	Für $n=3$ Gruppen mit Blockgrößen $32, 256, 64$ (aufsteigend in Dateireihenfolge) liefert die Sortierung die Reihenfolge $256, 64, 32$ und damit die Startadressen $0, 256, 320$. In aufsteigender Reihenfolge ($32, 64, 256$) lägen die Startadressen bei $0, 64, 256$ und der letzte Block endet erst bei $512$ statt bei $352$.
\end{beispiel}

\subsubsection{Firewall-Regeln aus der Zugriffsmatrix}

Alles, was nicht ausdrücklich in \code{access} freigegeben ist, wird verworfen (Default-Deny). Jedes Zugriffspaar $(i,j) \in P$ ergibt genau eine Freigabe vom Subnetz $N_i$ in das Subnetz $N_j$ (Satz~\ref{satz:firewall}).

